\subsection{\label{sec:mechanics-and-integration}Mechanics \& Design of Modern Detectors}

Modern detector systems are large, complicated systems. Integrating a Cherenkov detector (or any detector, for that matter) into such an system presents a unique set of mechanical and design challenges. Mechanical tolerances and alignment are crucial to maintaining accuracy. Optical elements must be precisely aligned and consistently calibrated throughout the detector lifetime. Careful thermal and radiation management are also necessary to ensure the longevity of photosensitive components. Consideration of the surrounding magnetic field from the solenoid used by the TOF detectors is also necessary, especially for PMTs which are magneto-sensitive. 

Ensuring that the detector can be serviced and maintained is also a key consideration, especially for large scale experiments that may run for decades. Some experiments are built on rails to allow the inner detectors (RICH, TOF, trackers etc.) the be slid out for maintenance. These considerations are imperative to the success of a detector.

RICH detectors require a large volume radially to generate and image the cone of Cherenkov photons. DIRC detectors on the other hand, are far more radially compact, offering an advantage in certain facilities. However, the detection of natural events is not possible with either of these systems. If one wanted to study the particles incident on Earth from outer-space, a different system entirely is needed.